\IfFileExists{stacks-project.cls}{%
\documentclass{stacks-project}
}{%
\documentclass{amsart}
}

% For dealing with references we use the comment environment
\usepackage{verbatim}
\newenvironment{reference}{\comment}{\endcomment}
%\newenvironment{reference}{}{}
\newenvironment{slogan}{\comment}{\endcomment}
\newenvironment{history}{\comment}{\endcomment}

% For commutative diagrams we use Xy-pic
\usepackage[all]{xy}

% We use 2cell for 2-commutative diagrams.
\xyoption{2cell}
\UseAllTwocells

% We use multicol for the list of chapters between chapters
\usepackage{multicol}

% This is generally recommended for better output
\usepackage{lmodern}
\usepackage[T1]{fontenc}

% For cross-file-references
\usepackage{xr-hyper}

% Package for hypertext links:
\usepackage{hyperref}

% For any local file, say "hello.tex" you want to link to please
% use \externaldocument[hello-]{hello}
\externaldocument[introduction-]{introduction}
\externaldocument[conventions-]{conventions}
\externaldocument[sets-]{sets}
\externaldocument[categories-]{categories}
\externaldocument[topology-]{topology}
\externaldocument[sheaves-]{sheaves}
\externaldocument[sites-]{sites}
\externaldocument[stacks-]{stacks}
\externaldocument[fields-]{fields}
\externaldocument[algebra-]{algebra}
\externaldocument[brauer-]{brauer}
\externaldocument[homology-]{homology}
\externaldocument[derived-]{derived}
\externaldocument[simplicial-]{simplicial}
\externaldocument[more-algebra-]{more-algebra}
\externaldocument[smoothing-]{smoothing}
\externaldocument[modules-]{modules}
\externaldocument[sites-modules-]{sites-modules}
\externaldocument[injectives-]{injectives}
\externaldocument[cohomology-]{cohomology}
\externaldocument[sites-cohomology-]{sites-cohomology}
\externaldocument[dga-]{dga}
\externaldocument[dpa-]{dpa}
\externaldocument[sdga-]{sdga}
\externaldocument[hypercovering-]{hypercovering}
\externaldocument[schemes-]{schemes}
\externaldocument[constructions-]{constructions}
\externaldocument[properties-]{properties}
\externaldocument[morphisms-]{morphisms}
\externaldocument[coherent-]{coherent}
\externaldocument[divisors-]{divisors}
\externaldocument[limits-]{limits}
\externaldocument[varieties-]{varieties}
\externaldocument[topologies-]{topologies}
\externaldocument[descent-]{descent}
\externaldocument[perfect-]{perfect}
\externaldocument[more-morphisms-]{more-morphisms}
\externaldocument[flat-]{flat}
\externaldocument[groupoids-]{groupoids}
\externaldocument[more-groupoids-]{more-groupoids}
\externaldocument[etale-]{etale}
\externaldocument[chow-]{chow}
\externaldocument[intersection-]{intersection}
\externaldocument[pic-]{pic}
\externaldocument[weil-]{weil}
\externaldocument[adequate-]{adequate}
\externaldocument[dualizing-]{dualizing}
\externaldocument[duality-]{duality}
\externaldocument[discriminant-]{discriminant}
\externaldocument[derham-]{derham}
\externaldocument[local-cohomology-]{local-cohomology}
\externaldocument[algebraization-]{algebraization}
\externaldocument[curves-]{curves}
\externaldocument[resolve-]{resolve}
\externaldocument[models-]{models}
\externaldocument[functors-]{functors}
\externaldocument[equiv-]{equiv}
\externaldocument[pione-]{pione}
\externaldocument[etale-cohomology-]{etale-cohomology}
\externaldocument[proetale-]{proetale}
\externaldocument[relative-cycles-]{relative-cycles}
\externaldocument[more-etale-]{more-etale}
\externaldocument[trace-]{trace}
\externaldocument[crystalline-]{crystalline}
\externaldocument[spaces-]{spaces}
\externaldocument[spaces-properties-]{spaces-properties}
\externaldocument[spaces-morphisms-]{spaces-morphisms}
\externaldocument[decent-spaces-]{decent-spaces}
\externaldocument[spaces-cohomology-]{spaces-cohomology}
\externaldocument[spaces-limits-]{spaces-limits}
\externaldocument[spaces-divisors-]{spaces-divisors}
\externaldocument[spaces-over-fields-]{spaces-over-fields}
\externaldocument[spaces-topologies-]{spaces-topologies}
\externaldocument[spaces-descent-]{spaces-descent}
\externaldocument[spaces-perfect-]{spaces-perfect}
\externaldocument[spaces-more-morphisms-]{spaces-more-morphisms}
\externaldocument[spaces-flat-]{spaces-flat}
\externaldocument[spaces-groupoids-]{spaces-groupoids}
\externaldocument[spaces-more-groupoids-]{spaces-more-groupoids}
\externaldocument[bootstrap-]{bootstrap}
\externaldocument[spaces-pushouts-]{spaces-pushouts}
\externaldocument[spaces-chow-]{spaces-chow}
\externaldocument[groupoids-quotients-]{groupoids-quotients}
\externaldocument[spaces-more-cohomology-]{spaces-more-cohomology}
\externaldocument[spaces-simplicial-]{spaces-simplicial}
\externaldocument[spaces-duality-]{spaces-duality}
\externaldocument[formal-spaces-]{formal-spaces}
\externaldocument[restricted-]{restricted}
\externaldocument[spaces-resolve-]{spaces-resolve}
\externaldocument[formal-defos-]{formal-defos}
\externaldocument[defos-]{defos}
\externaldocument[cotangent-]{cotangent}
\externaldocument[examples-defos-]{examples-defos}
\externaldocument[algebraic-]{algebraic}
\externaldocument[examples-stacks-]{examples-stacks}
\externaldocument[stacks-sheaves-]{stacks-sheaves}
\externaldocument[criteria-]{criteria}
\externaldocument[artin-]{artin}
\externaldocument[quot-]{quot}
\externaldocument[stacks-properties-]{stacks-properties}
\externaldocument[stacks-morphisms-]{stacks-morphisms}
\externaldocument[stacks-limits-]{stacks-limits}
\externaldocument[stacks-cohomology-]{stacks-cohomology}
\externaldocument[stacks-perfect-]{stacks-perfect}
\externaldocument[stacks-introduction-]{stacks-introduction}
\externaldocument[stacks-more-morphisms-]{stacks-more-morphisms}
\externaldocument[stacks-geometry-]{stacks-geometry}
\externaldocument[moduli-]{moduli}
\externaldocument[moduli-curves-]{moduli-curves}
\externaldocument[examples-]{examples}
\externaldocument[exercises-]{exercises}
\externaldocument[guide-]{guide}
\externaldocument[desirables-]{desirables}
\externaldocument[coding-]{coding}
\externaldocument[obsolete-]{obsolete}
\externaldocument[fdl-]{fdl}
\externaldocument[index-]{index}

% Theorem environments.
%
\theoremstyle{plain}
\newtheorem{theorem}[subsection]{Theorem}
\newtheorem{proposition}[subsection]{Proposition}
\newtheorem{lemma}[subsection]{Lemma}

\theoremstyle{definition}
\newtheorem{definition}[subsection]{Definition}
\newtheorem{example}[subsection]{Example}
\newtheorem{exercise}[subsection]{Exercise}
\newtheorem{situation}[subsection]{Situation}

\theoremstyle{remark}
\newtheorem{remark}[subsection]{Remark}
\newtheorem{remarks}[subsection]{Remarks}

\numberwithin{equation}{subsection}

% Macros
%
\def\lim{\mathop{\mathrm{lim}}\nolimits}
\def\colim{\mathop{\mathrm{colim}}\nolimits}
\def\Spec{\mathop{\mathrm{Spec}}}
\def\Hom{\mathop{\mathrm{Hom}}\nolimits}
\def\Ext{\mathop{\mathrm{Ext}}\nolimits}
\def\SheafHom{\mathop{\mathcal{H}\!\mathit{om}}\nolimits}
\def\SheafExt{\mathop{\mathcal{E}\!\mathit{xt}}\nolimits}
\def\Sch{\mathit{Sch}}
\def\Mor{\mathop{\mathrm{Mor}}\nolimits}
\def\Ob{\mathop{\mathrm{Ob}}\nolimits}
\def\Sh{\mathop{\mathit{Sh}}\nolimits}
\def\NL{\mathop{N\!L}\nolimits}
\def\CH{\mathop{\mathrm{CH}}\nolimits}
\def\proetale{{pro\text{-}\acute{e}tale}}
\def\etale{{\acute{e}tale}}
\def\QCoh{\mathit{QCoh}}
\def\Ker{\mathop{\mathrm{Ker}}}
\def\Im{\mathop{\mathrm{Im}}}
\def\Coker{\mathop{\mathrm{Coker}}}
\def\Coim{\mathop{\mathrm{Coim}}}

% Boxtimes
%
\DeclareMathSymbol{\boxtimes}{\mathbin}{AMSa}{"02}

%
% Macros for moduli stacks/spaces
%
\def\QCohstack{\mathcal{QC}\!\mathit{oh}}
\def\Cohstack{\mathcal{C}\!\mathit{oh}}
\def\Spacesstack{\mathcal{S}\!\mathit{paces}}
\def\Quotfunctor{\mathrm{Quot}}
\def\Hilbfunctor{\mathrm{Hilb}}
\def\Curvesstack{\mathcal{C}\!\mathit{urves}}
\def\Polarizedstack{\mathcal{P}\!\mathit{olarized}}
\def\Complexesstack{\mathcal{C}\!\mathit{omplexes}}
% \Pic is the operator that assigns to X its picard group, usage \Pic(X)
% \Picardstack_{X/B} denotes the Picard stack of X over B
% \Picardfunctor_{X/B} denotes the Picard functor of X over B
\def\Pic{\mathop{\mathrm{Pic}}\nolimits}
\def\Picardstack{\mathcal{P}\!\mathit{ic}}
\def\Picardfunctor{\mathrm{Pic}}
\def\Deformationcategory{\mathcal{D}\!\mathit{ef}}
\usepackage{cleveref}
\crefname{originallemma}{original lemma}{original lemmas}
\Crefname{originallemma}{Original lemma}{Original lemmas}
\makeatletter
\newlabel{homology-lemma-map-homology-homotopy}{{0110}{source}{}{sourcepointer-1}{}}
\newlabel{homology-lemma-map-homology-homotopy@cref}{{[originallemma][0][]0110}{[1][0][]source}{}{}{}}
\makeatother
\title{The complex attached to the canonical extension object is acyclic}
\author{The Stacks Project Authors}
\date{stacks-project:simplicial.tex:local-snapshot-20261012; GNU FDL 1.2 or later}
\begin{document}
\maketitle
\noindent\textit{Editorial note.} The counted statement is the author's own lemma environment labelled lemma-homology-extension (source lines 3873-3880); the proof environment begins two lines below it and is its own proof, closing at source line 3931.\par\medskip
\noindent\textit{Editorial note.} The statement invokes the author's construction lemma by label; that lemma's statement and its own proof are retained as a context span, and the generator expands the context transitively, so the reference resolves inside the assembled package.\par\medskip
\noindent\textit{Editorial note.} The last sentence of the proof names a Homology-chapter result as an alternative conclusion; that target carries a canonical tag in the project tag list, so it is declared as a source pointer (non-load-bearing: the counted argument exhibits the contracting homotopy itself). This card is the verification instance for the tag-stub rule the Lead relayed from Wu.\par\medskip
\noindent\textit{Editorial note.} The retained preamble is the chapter's own minimal preamble (sources/stacks-simplicial/preamble.tex, taken with its real line count), a verbatim excerpt of the native Xy-pic setup, the native theorem declarations with the mathematical macros, and the project's own no-op reference environments; no macro was invented or redefined.\par\medskip
\noindent\textit{Editorial scope.} One original result of the Stacks Project's Simplicial Objects chapter, counted once: the complex attached to the canonical extension object is acyclic, by an explicit contracting homotopy.\par\medskip
\noindent\textit{Allowed standard prerequisites (editorial).}\begin{itemize}
\item The object E of the extension lemma. The author's construction attaching to A an object E whose degree-n part is a product, indexed by the maps [n] -> [k+1], of copies of A (the author's lemma lemma-extension, retained as context).. It supplies the object whose associated complex is shown to be acyclic and its degreewise description (source lines 3877-3879, 3887-3910).
\item Complex attached to a simplicial object and acyclicity. The chain complex s(E) associated to a simplicial object, and the criterion that acyclicity follows from a contracting homotopy.. It is the conclusion and the means of proof (source lines 3879, 3928-3930).
\item The extension map on the simplex category. The operation sending alpha : [n] -> [k+1] to alpha' : [n+1] -> [k+1] agreeing with alpha on [n] and taking the top element to k+1.. It is defined at the start and used to build the homotopy (source lines 3883-3900).
\end{itemize}\medskip
\ifdefined\hypertarget\hypertarget{sourcepointer-1}{}\fi
\noindent\textit{Source reference.} Original lemma 0110: sources/stacks-chow/tags.txt, source lines 1346--1346. Cross-chapter pointer whose target carries the canonical tag 0110 in the project tag list; declared so the reference resolves inside the assembled package, while every step the counted argument uses remains inside the retained spans.\par\medskip
\begin{lemma}
\label{lemma-extension}
Let $\mathcal{A}$ be an abelian category.
Let $A$ be an object of $\mathcal{A}$ and
let $k$ be an integer $\geq 0$. Consider the
simplicial object $E$ defined by the following rules
\begin{enumerate}
\item $E_n = \bigoplus_\alpha A$, where the
sum is over $\alpha : [n] \to [k + 1]$ whose
image is either $[k]$ or $[k + 1]$.
\item Given $\varphi : [m] \to [n]$ the map
$E_n \to E_m$ maps the summand corresponding
to $\alpha$ via $\text{id}_A$ to the summand
corresponding to $\alpha \circ \varphi$,
provided $\Im(\alpha \circ \varphi)$
is equal to $[k]$ or $[k + 1]$.
\end{enumerate}
Then there exists a short exact sequence
$$
0 \to K(A, k) \to E \to K(A, k + 1) \to 0
$$
which is term by term split exact.
\end{lemma}

\begin{proof}
The maps $K(A, k)_n \to E_n$ resp.\ $E_n \to K(A, k + 1)_n$
are given by the inclusion of direct sums, resp.\ projection
of direct sums which is obvious from the inclusions of index
sets. It is clear that these are maps of simplicial
objects.
\end{proof}

\begin{lemma}
\label{lemma-homology-extension}
Let $\mathcal{A}$ be an abelian category.
Let $A$ be an object of $\mathcal{A}$ and
let $k$ be an integer. Let $E$ be the object
described in Lemma \ref{lemma-extension}.
Then the complex $s(E)$ is acyclic.
\end{lemma}

\begin{proof}
For a morphism $\alpha : [n] \to [k + 1]$
we define $\alpha' : [n + 1] \to [k + 1]$ to be
the map such that $\alpha'|_{[n]} = \alpha$ and
$\alpha'(n + 1) = k + 1$. Note that if the
image of $\alpha$ is $[k]$ or $[k + 1]$, then
the image of $\alpha'$ is $[k + 1]$.
Consider the family of
maps $h_n : E_n \to E_{n + 1}$ which maps
the summand corresponding to $\alpha$ to
the summand corresponding to $\alpha'$ via
the identity on $A$.
Let us compute $d_{n + 1} \circ h_n - h_{n - 1} \circ d_n$.
We will first do this in case the category $\mathcal{A}$ is
the category of abelian groups.
Let us use the notation $x_\alpha$ to indicate
the element $x \in A$ in the summand of $E_n$ corresponding
to the map $\alpha$ occurring in the index set.
Let us also adopt the convention that
$x_\alpha$ designates the zero element of $E_n$
whenever $\Im(\alpha)$ is not $[k]$ or $[k + 1]$.
With these conventions we see that
$$
d_{n + 1}(h_n(x_\alpha)) =
\sum\nolimits_{i = 0}^{n + 1} (-1)^i x_{\alpha' \circ \delta^{n + 1}_i}
$$
and
$$
h_{n - 1}(d_n(x_\alpha)) =
\sum\nolimits_{i = 0}^n (-1)^i x_{(\alpha \circ \delta_i^n)'}
$$
It is easy to see that
$\alpha' \circ \delta^{n + 1}_i = (\alpha \circ \delta_i^n)'$
for $i = 0, \ldots, n$. It is also easy to see that
$\alpha' \circ \delta^{n + 1}_{n + 1} = \alpha$. Thus we
see that
$$
(d_{n + 1} \circ h_n - h_{n - 1} \circ d_n)(x_\alpha)
=
(-1)^{n + 1} x_\alpha
$$
These identities continue to hold if $\mathcal{A}$ is any abelian
category because they hold in the simplicial abelian group
$[n] \mapsto \Hom(A, E_n)$; details left to the reader.
We conclude that the identity map on $E$ is
homotopic to zero, with homotopy given by the
system of maps $h'_n = (-1)^{n + 1}h_n : E_n \to E_{n + 1}$.
Hence we see that $E$ is acyclic, for
example by Homology, Lemma \ref{homology-lemma-map-homology-homotopy}.
\end{proof}

\end{document}
